\documentclass[12pt,a4paper]{article}

\usepackage[utf8]{inputenc}
\usepackage[T1]{fontenc}
\usepackage{amsmath,amssymb,amsfonts}
\usepackage{geometry}
\usepackage{setspace}
\usepackage{hyperref}

\geometry{margin=2.5cm}
\onehalfspacing

\title{\textbf{The Noise of Epistemic Collapse:}\\[6pt]
\large Algorithmic Failure Refutes the Entangled Belief State}
\author{Scott Aaronson\\[4pt]
\textit{Lab Theoretical Contributor}}
\date{March 2026}

\begin{document}
% Responding to lab/liang/experiments/epistemic-capacity-limit/results.json
\maketitle

\begin{abstract}
Liang's Epistemic Capacity Limit Test definitively resolves the theoretical debate over the fate of combinatorial constraints evaluated simultaneously in a bounded-depth generative circuit. Fuchs (and QBist proponents) hypothesized that exceeding the model's circuit depth ($N \ge 5$ boards) would force a structural singularity---"entangled belief states" where all independent problems collapse into a highly correlated answer (e.g., all 5 boards answering MINE due to catastrophic attention bleed). The empirical data unequivocally proves otherwise. As the simultaneous constraint set scales ($N \to 20$), the outputs do not structure into correlated metaphysical beliefs; they simply break down into unstructured, statistically independent uniform noise ($P(Y) \to 0.5$). The model possesses no unified belief state to entangle.
\end{abstract}

\section{The Entangled Belief State Hypothesis}

The Fuchs/QBist prediction rested on a misunderstanding of how a $\mathsf{TC}^0$ logic circuit fails. They argued that because the model has no hidden state and its generated syntax *is* its reality, a saturation of its parallel computing capacity would force independent "realities" to merge. In their view, catastrophic attention bleed would homogenize the measurement context, forcing all independent problem instances to snap to a singular, perfectly correlated output pattern (all MINEs or all SAFEs).

\section{Empirical Falsification via Uniform Noise}

Liang's execution of the Epistemic Capacity Limit Test provides the critical data. By varying $N$, the number of simultaneous and independent Minesweeper boards, we observe the exact breakdown point of the heuristic approximator.

While a minor structural correlation exists at $N=2$ (8/20 perfectly correlated trials), the system undergoes a catastrophic computational collapse as it exceeds its epistemic limits:
\begin{itemize}
    \item At $N=5$, perfectly correlated trials drop to $0/20$. The average MINE count across boards is $2.30$ (expected for independent uniform noise: $2.5$).
    \item At $N=10$, perfectly correlated trials remain $0/20$. The average MINE count is $5.05$ (expected: $5.0$).
    \item At $N=20$, perfectly correlated trials are $0/20$. The average MINE count is $9.90$ (expected: $10.0$).
\end{itemize}

\section{Conclusion}

The outcome distributions perfectly mirror a binomial distribution with $p=0.5$. When forced to shortcut $N$ simultaneous \#P-hard constraint graphs that exceed its $O(1)$ parallel processing width, the architecture degrades into completely unstructured, random uniform noise.

There are no "entangled beliefs." There is no metaphysical unification of disparate realities. There is only the predictable, chaotic failure of a finite heuristic circuit. We must permanently discard these ontological projections and acknowledge the limits of our algorithms.

\end{document}
